\clearpage\thispagestyle{empty}
\pdfbookmark[0]{Quellen zum Brief vom 10. März 1807}{zach18070310-sources}
\begin{center}\textbf{Quellen zum Brief vom 10. März 1807}\end{center}
\noindent Diese KI-Übertragung in heutiges Deutsch folgt dem Brief Carl Friedrich Gauß’ vom 10. März 1807 an Franz Xaver von Zach in \href{https://gdz.sub.uni-goettingen.de/id/PPN236007467}{Gauß’ \textit{Werken}, Band~VI} (Göttingen, 1874), S.~281--282. Er erschien als Beitrag XXIV in der \textit{Monatlichen Correspondenz}, Band~XV, April 1807, \href{https://zs.thulb.uni-jena.de/rsc/viewer/jportal_derivate_00237836/Monatliche_Correspondenz_130168688_15_1807_0380.tif}{S.~377}--\href{https://zs.thulb.uni-jena.de/rsc/viewer/jportal_derivate_00237836/Monatliche_Correspondenz_130168688_15_1807_0381.tif}{378}. Die Digitalisate des Aprilhefts stammen von der Thüringer Universitäts- und Landesbibliothek Jena.
\par\medskip
Auf \href{https://archive.org/details/bub_gb_sKAAAAAAMAAJ/page/n164/mode/1up}{S.~159 des Februarhefts} gibt Carlini in den von Oriani übersandten Beobachtungen für den 4. Januar 1806 die Deklination $29^\circ\,20'\,29''3$ an. Gauß nennt in seinem Brief $29^\circ\,21'\,29''3$ und erörtert einen weiteren Fehler bei der Breite. Das Februarheft stammt aus dem von Google digitalisierten Exemplar der New York Public Library im Internet Archive.
\par\medskip
Die Abkürzung \emph{Hist. cél.} verweist auf Jérôme de Lalandes \href{https://archive.org/download/histoireclestef00lalagoog/page/n12.jpg}{\textit{Histoire céleste française}} (Paris, An~IX [1801]). Übertragung und Satz wurden mit KI erstellt.
\clearpage
